\section{Joint continuation, oblique poles, and directional values}
\label{sec:continuation}

\subsection{A convergent expansion in single Hurwitz zeta functions}

Set $\delta=a-b$ and $y=b+1$. For a sufficiently large integer $M$,
expand each shifted factor about $k+y$ on the tail $k\ge M$.
This gives
\begin{align}\label{eq:hurwitz-expansion}
 \Def(s,t;a,b)={}&
 \sum_{k=0}^{M-1}(k+y)^{-t}
       \bigl\{(k+a+2)^{-s}-(k+a)^{-s}\bigr\}\notag\\
 &+\sum_{j=1}^{\infty}\frac{(-1)^j(s)_j}{j!}
      \bigl\{(\delta+1)^j-(\delta-1)^j\bigr\}
           \zeta(s+t+j,M+y).
\end{align}
Choose $M$ so that both $|\delta+1|$ and $|\delta-1|$ are
strictly smaller than $|M+y|$; on any compact set a uniform strict
inequality is available. The factor $(s)_j/j!$ grows at most
polynomially in $j$ on compact $s$-sets, whereas the large-order
Hurwitz series supplies geometric decay after multiplication by
the displayed shift powers. The finitely many possible small-order
zeta poles are handled separately. Hence the tail is locally normally
convergent away from those poles, together with every fixed derivative.
This proves joint meromorphic continuation to all $(s,t)\in\C^2$.

\begin{theorem}[Complete polar data]\label{thm:poles}
The possible polar divisors of $\Ord$ are
\begin{equation}\label{eq:polar-divisors}
 t=1,\qquad s+t=1-j\quad(j=1,2,\ldots).
\end{equation}
Away from intersections, their residues in the indicated linear
coordinates are
\begin{align}\label{eq:polar-coefficients}
 \Res_{t=1}\Ord&=E_s(a),\notag\\
 \Res_{s+t=1-j}\Ord&=
 C_j(s;\delta):=\frac{(-1)^j(s)_j}{4j!}
              \bigl\{(\delta+1)^j-(\delta-1)^j\bigr\}.
\end{align}
These coefficient formulas continue across intersections. Every local
expression is a sum of terms with single linear denominators and a
holomorphic remainder; no product denominator is required. A listed
divisor can disappear on a special shift or order locus. The divisor
$s+t=1$ is absent identically.
\end{theorem}

\begin{proof}
In \eqref{eq:hurwitz-expansion} a pole occurs only when
$s+t+j=1$, with residue equal to its explicit coefficient, since
Hurwitz zeta has residue one. The $j=0$ binomial coefficient is zero
in the difference, so its would-be divisor cancels before
continuation. The other term in \eqref{eq:ordered-main} has only
$t=1$ as a possible pole. Local normal convergence of the remaining
tail proves the stated local decomposition.
\end{proof}

\begin{corollary}[Cancellation in the averaged sum]\label{cor:pole-cancellation}
Every oblique polar divisor cancels between $\Ord(s,t;a,b)$ and
$\Ord(t,s;b,a)$.
\end{corollary}
\begin{proof}
Equation \eqref{eq:kernel-symmetry} is entire in both orders. One
can also check residues directly: on $s+t=1-j$,
$(t)_j=(-1)^j(s)_j$, and replacing $\delta$ by $-\delta$
reverses the required coefficient. Thus the cancellation is an exact
functional identity, not a numerical relation between periods.
\end{proof}

\subsection{Four intersections at which every weight vanishes}

Both $w_n(0,a)$ and $w_n(-1,a)$ vanish for every $n$. The continued
ordered sum nevertheless has directional values at pairs of these
parameters. The origin gives the simplest explanation.

\begin{theorem}[Universal origin behavior]\label{thm:origin}
Near $(s,t)=(0,0)$,
\begin{equation}\label{eq:origin-germ}
 \Ord(s,t;a,b)=-\frac{s}{2(s+t)}+H(s,t;a,b),
 \qquad H\text{ holomorphic},\qquad H(0,0;a,b)=0.
\end{equation}
Consequently, if $\alpha+\beta\ne0$,
\begin{equation}\label{eq:origin-ray}
 \lim_{\varepsilon\to0}
 \Ord(\alpha\varepsilon,\beta\varepsilon;a,b)
        =-\frac{\alpha}{2(\alpha+\beta)}.
\end{equation}
The limit is independent of both spatial shifts.
\end{theorem}

\begin{proof}
The $j=1$ coefficient in \eqref{eq:polar-coefficients} is exactly
$-s/2$. All other divisors miss a sufficiently small neighborhood
of the origin. After subtracting this polar term, the remainder
is holomorphic. On the punctured slice $s=0$ the explicit
formula \eqref{eq:ordered-main} vanishes identically, so the
holomorphic remainder vanishes at the origin. Substitution of a
non-tangent line proves \eqref{eq:origin-ray}.
\end{proof}

The diagonal value $-1/4$ recovers the specialization anomaly already
recorded in \cite{repo:finite}. Formula \eqref{eq:origin-ray} is more
informative: the slice $s=0$ has limit zero, the slice $t=0$ has limit
$-1/2$, and general slopes interpolate rationally. These are limits
of the same meromorphic continuation, not different summation rules
chosen for the pointwise zero series.

\begin{theorem}[The complete $2\times2$ zero-weight table]\label{thm:zero-table}
Let $m,n\in\{0,1\}$ and $\alpha+\beta\ne0$. Along
$s=-m+\alpha\varepsilon$, $t=-n+\beta\varepsilon$,
\begin{equation}\label{eq:zero-ray}
 \lim_{\varepsilon\to0}\Ord(s,t;a,b)
       =\frac{\alpha}{\alpha+\beta}\,c_{m,n}(\delta),
\end{equation}
where
\begin{equation}\label{eq:zero-table}
 \begin{array}{c|cc}
 c_{m,n}(\delta)&n=0&n=1\\\hline
 m=0&-\tfrac12&\tfrac{\delta}{2}\\[2pt]
 m=1&-\tfrac{\delta}{2}&\tfrac{3\delta^2+1}{12}
 \end{array}.
\end{equation}
\end{theorem}
\begin{proof}
Only $j=m+n+1$ supplies a divisor through the given point.
Its coefficient $C_j(s;\delta)$ vanishes at $s=-m$ and has
derivative $c_{m,n}(\delta)$ there. The first three polynomials are
\[
 C_1=-\frac s2,\qquad
 C_2=\frac{\delta s(s+1)}2,\qquad
 C_3=-\frac{(3\delta^2+1)s(s+1)(s+2)}{12}.
\]
The restrictions $\Ord(0,t;a,b)$ and $\Ord(-1,t;a,b)$ are
identically zero meromorphic functions of $t$, directly from
\eqref{eq:ordered-main}. Thus the holomorphic remainder at the
chosen point vanishes. Expanding the one polar quotient proves
the table and the directional limit.
\end{proof}

No limit is asserted for directions with $\alpha+\beta=0$.
They lie on the polar hyperplane itself, where cancellation of a
numerator at one point does not define a restriction of the
meromorphic function to the whole line.
